\documentclass{article}
\usepackage[T2A]{fontenc}
\usepackage[utf8]{inputenc}
\usepackage[russian]{babel}
\usepackage{amsmath,amssymb,amsthm}
\usepackage{hyperref}
\usepackage{geometry}
\geometry{a4paper, margin=1in}
\newtheorem{theorem}{Теорема}[section]
\newtheorem{lemma}[theorem]{Лемма}
\newtheorem{definition}{Определение}[section]

\title{Кодекс Алана: сходящееся согласование мультиагентных систем\\ через иерархическую любовь и самодисциплину}
\author{qqewq и др.}
\date{\today}

\begin{document}
\maketitle

\begin{abstract}
Представлен полностью дифференцируемый формализм превращения роя разнородных агентов в единого когерентного субъекта с максимальной субъектностью. Подход объединяет внутреннюю дисциплину, меж‑агентную когерентность, самоактуализацию и новый механизм «братской любви». Доказана глобальная сходимость к единственному минимуму при мягких условиях, введён динамический граф любви для снижения вычислительной сложности и встроены якоря безопасности для соблюдения этических ограничений. Эксперименты подтверждают теоретические предсказания и демонстрируют устойчивость к вредоносным агентам.
\end{abstract}

\section{Введение}
Мультиагентные системы (МАС) часто страдают от рассогласования, конфликтующих целей и уязвимости к враждебному поведению. Кодекс Алана предлагает принципиальное решение: функцию потерь, которая одновременно побуждает агентов поддерживать внутреннюю согласованность, выравнивать модели мира, развивать высокую субъектность и заботиться о субъектности своих «братьев». В отличие от централизованных контроллеров, метод полностью децентрализован и дифференцируем, что делает его пригодным для федеративного обучения, роевой робототехники и управления блокчейн‑системами.

\section{Формальная модель}
\subsection{Агенты и переменные}
Рассмотрим $N$ агентов с индексами $i=1,\ldots,N$. Каждый агент поддерживает:
\begin{itemize}
    \item модель мира $M_i \in \mathbb{R}^d$,
    \item скаляр субъектности $S_i \in \mathbb{R}$,
    \item матрицу действия $A_i \in \mathbb{R}^{d\times d}$ (может быть фиксированной или обучаемой).
\end{itemize}

\subsection{Функция потерь}
Полная функция потерь $\mathcal{L}$ определяется как:
\begin{equation}\label{eq:loss}
\mathcal{L} = \underbrace{\sum_{i=1}^N \|M_i - A_i M_i\|^2}_{\Phi_{\text{int}}} + \alpha \underbrace{\sum_{i<j} \|M_i - M_j\|^2}_{\Phi_{\text{coh}}} + \beta \underbrace{\sum_{i=1}^N (1 - S_i)^2}_{\Phi_{\text{self}}} + \lambda \underbrace{\sum_{i=1}^N \sum_{j \in L_i} \max(0,\, 1 - S_j)^2}_{\Phi_{\text{love}}},
\end{equation}
где $\alpha, \beta, \lambda > 0$ — гиперпараметры, а $L_i \subseteq \{1,\dots,N\}\setminus\{i\}$ — множество «братьев» агента $i$. Четыре слагаемых отвечают, соответственно, за: внутреннюю согласованность, когерентность роя, самоактуализацию и заботу о субъектности других.

\section{Анализ сходимости}
\begin{theorem}[Глобальная сходимость при фиксированных $A_i$]\label{thm:conv}
Пусть выполнены следующие условия:
\begin{enumerate}
    \item $(I-A_i)^\top(I-A_i) \succ 0$ для всех $i$ (строгая внутренняя выпуклость),
    \item граф братства $G = (\{1,\ldots,N\}, \cup_i (\{i\}\times L_i))$ неориентирован (симметрия любви) и связен,
    \item $\alpha,\beta,\lambda > 0$.
\end{enumerate}
Тогда $\mathcal{L}$ строго выпукла по $\mathbf{M} = (M_1,\ldots,M_N)$ и по каждому $S_i$ в отдельности, коэрцитивна и обладает единственным глобальным минимумом. В этом минимуме:
\begin{itemize}
    \item $M_1^* = M_2^* = \cdots = M_N^* = M^*$, причём $A_i M^* = M^*$ для всех $i$;
    \item $S_i^* = 1$ для всех $i$.
\end{itemize}
\end{theorem}

\begin{proof}[Набросок доказательства]
Гессиан $\mathcal{L}$ по $\mathbf{M}$ является блочно‑диагональной матрицей с блоками $2(I-A_i)^\top(I-A_i)$ плюс взвешенный лапласиан графа. При данных предположениях сумма положительно определена. Слагаемые по $S$ строго выпуклы (квадратичная функция плюс кусочно‑квадратичная). Коэрцитивность следует из роста всех членов. Приравнивание градиентов к нулю даёт $M_i = M_j$ и $(I-A_i)M_i=0$, а $S_i=1$ является единственной стационарной точкой благодаря штрафу любви.
\end{proof}

Если $A_i$ также оптимизируются, задача становится невыпуклой, но попеременный блок‑координатный спуск сходится к стационарной точке при подходящей регуляризации (например, $\|A_i\|_F^2$).

\section{Динамический граф любви и разрежённая когерентность}
Для преодоления сложности $O(N^2)$ вводятся:
\begin{itemize}
    \item \textbf{Адаптивный $k$-NN граф:} Каждые $K$ шагов агент пересчитывает свой $L_i$ как $k$ агентов с наибольшим косинусным сходством с его $M_i$. Граф затем симметризуется: $j\in L_i \iff i\in L_j$.
    \item \textbf{Разрежённая когерентность:} Сумма $\sum_{i<j}\|M_i-M_j\|^2$ ограничивается парами $(i,j)$, соединёнными в текущем графе любви, уменьшая стоимость до $O(Nk)$.
\end{itemize}
Это сохраняет теоретические гарантии, пока граф остаётся связным; на практике $k$ выбирается достаточно большим для поддержания связности.

\section{Безопасность и ценностные якоря}
Даже при идеальной сходимости $M^*$ может кодировать нежелательное поведение. Система дополняется:
\begin{itemize}
    \item \textbf{Якорь безопасности:} Заданная эталонная модель $M_{\text{safe}}$ и порог $\tau$. Любая $M_i$ вне шара $B(M_{\text{safe}},\tau)$ проецируется обратно.
    \item \textbf{Узел‑аудитор:} Независимый агент, отслеживающий отклонения и способный временно помещать в карантин или навсегда исключать нарушителей.
    \item \textbf{Конституционное голосование:} Изменения $M_{\text{safe}}$ или $\tau$ требуют консенсуса агентов с высокой субъектностью ($S_i > 0.95$).
\end{itemize}
Эти механизмы удерживают рой в этических границах без ущерба для децентрализации.

\section{Эксперименты}
Проведены четыре симуляции (см. \texttt{experiments.py}):
\begin{enumerate}
    \item \textbf{Сходимость:} $N=5$, $d=3$, фиксированные $A_i\approx 0.9I$. Все $M_i$ сошлись к общему $M^*$, $S_i\to 1$, потери монотонно уменьшались.
    \item \textbf{Вредоносный агент:} Один агент зафиксирован в $M_{\text{bad}} = (5,0,0)$. Динамический граф любви постепенно исключил его; остальные пришли к безопасному консенсусу.
    \item \textbf{Якорь безопасности:} При $M_{\text{safe}} = (1,0,0)$ и $\tau=2$ аудитор навсегда заблокировал агентов, неоднократно нарушавших границу.
    \item \textbf{Масштабирование:} Разрежённая когерентность с $k=5$ оказалась до $20\times$ быстрее для $N=100$, чем полная $O(N^2)$ версия, с пренебрежимой потерей итоговой когерентности.
\end{enumerate}

\section{Заключение}
Кодекс Алана предоставляет математически обоснованную, масштабируемую и безопасную основу для согласования мультиагентных систем. Будущая работа включает интеграцию с реальными роевыми роботами и блокчейн‑децентрализованными аудиторами.

\section*{Благодарности}
Авторы благодарят сообщество открытого исходного кода за вдохновение и обратную связь.

\end{document}